Enlarging the measured region raises body-glare recovery from 8/12 to
12/12 and removes all 18 raw errors (Table~\ref{tab:frontends}). Every
additional method also reaches 12/12. The improvement is therefore already
explained by measuring a larger area; these recordings show no further
benefit from filtering or reference compensation. In calibration, spatial
and combined processing recovered 2/2 glare packets where the larger mean
recovered 0/2. That motivated testing, but its advantage did not carry over
to the main glare comparison.

Under blockage, the larger mean raises recovery from 27/42 to 29/42,
rescuing three messages while losing one previously recovered. EWMA gives
30/42 and Kalman 29/42. Reference and combined processing reach 33/42,
each rescuing six messages without losing any baseline successes.
However, some blockages cover the reference regions, so these improvements
do not establish that compensation correctly estimated illumination changes.
EWMA has the fewest raw errors (71), yet recovers fewer messages than the
reference method (72 errors). A lower error count therefore does not
guarantee better recovery: correction also depends on how errors are grouped.

The benefit also depends on the code. Combined processing raises uncoded
blockage recovery from 1/6 to 2/6 and non-interleaved Hamming from 4/6 to
5/6. Enlarging the region alone instead lowers Hamming to 3/6 despite
improving the overall total. Interleaved Hamming and non-interleaved
Reed--Solomon remain at 6/6 throughout. A receiver improvement in aggregate
therefore need not improve every coding configuration.

These results favour measuring a sufficiently broad screen area before
adding more processing. Temporal filtering offers little additional benefit
here, while reference compensation depends on suitable reference signals.
All seven methods retain the six early-glare failures because they share
packet detection; this experiment cannot establish improved acquisition.
